%%%%%%%%%%%%%%%%%%%%%%%%%%%%%%%%%%%%%%%%%
% Arsclassica Article
% Structure Specification File
% Adapted for this project with XeLaTeX and Chinese content.
%%%%%%%%%%%%%%%%%%%%%%%%%%%%%%%%%%%%%%%%%

\usepackage[
nochapters,
beramono,
pdfspacing,
dottedtoc
]{classicthesis}

\usepackage{arsclassica}

\usepackage{fontspec}
\usepackage{xeCJK}
\usepackage{unicode-math}

\setmainfont{TeX Gyre Pagella}
\setsansfont{TeX Gyre Heros}
\setmonofont{Consolas}
\setCJKmainfont{SimSun}[BoldFont=SimHei]
\setCJKsansfont{SimHei}[BoldFont=SimHei]
\setCJKmonofont{KaiTi}
\setmathfont{TeX Gyre Pagella Math}

\usepackage{graphicx}
\graphicspath{{figures/}{../output/figures/}}

\usepackage{enumitem}
\usepackage{indentfirst}
\usepackage{subfig}
\usepackage{amsthm}
\usepackage{booktabs}
\usepackage{xpatch}
\usepackage{varioref}
\usepackage{bookmark}

\setlength{\parindent}{2em}
\setlist{itemsep=0.2em, topsep=0.3em}

\renewcommand{\contentsname}{目录}
\renewcommand{\listfigurename}{插图目录}
\renewcommand{\listtablename}{表格目录}
\renewcommand{\abstractname}{摘要}
\renewcommand{\refname}{参考文献}

\theoremstyle{definition}
\newtheorem{definition}{定义}[section]

\theoremstyle{plain}
\newtheorem{theorem}{定理}[section]
\newtheorem{lemma}[theorem]{引理}
\newtheorem{corollary}[theorem]{推论}

\theoremstyle{remark}


\hypersetup{
colorlinks=true,
breaklinks=true,
bookmarks=true,
bookmarksnumbered=true,
urlcolor=webbrown,
linkcolor=RoyalBlue,
citecolor=webgreen
}
